Scanning is the first step of the pipeline. Here we translate the character sequence from the source program into a token sequence.
We do this by iterating the entire input string trying to recognize different types of tokens.
%! Cuttet fra for plads. Er nok heller ikke så vigtigt?
%Each Token object can store the lexeme, literal, the line the token appears at and the type of the token which is one of the values from enum \texttt{TokenType}.
%To track where in the file it is at, the Scanner uses three pointers: start, current and line.
%The start pointer always indicates the start of the lexeme being scanned while current points to the character
%we are currently inspecting. The line field keeps track of the line we are at in the source program.
In each iteration we scan a single token by utilizing the \texttt{scanToken} method consiting of
one big switch statement to inspect a character and match it to the corresponding case.

The switch statement has a series of simple single character cases which
includes parenthesies, braces, coomas, semicolons, and minuses.
All of these cases are easy - we simply generate the a token with the corresponding type.
Other single character lexemes include comments, newlines and whitespaces.
All of these aren't meaningful to the parser and as such we perform the necessary logic to skip them. 


When encountering longer lexemes like strings, numbers and identifiers we call separate functions to handle these.
For strings we attempt to consume until the closing " is found reporting an error if the string spans too many lines.
When we see a number, we consume the entire series of digits allowing for a decimal point with one or more trailing digits.
%! Kan måske godt bare cuttes til: "leading and trailing decimal points are not allowed".
By looking ahead two characters we ensure a digit is present after the '\texttt{.}' disallowing trailing decimal points.
Leading decimal points are also not allowed, as the character '\texttt{.}' matches to no case in the switch statement.
The literal value is stored as a Double.

Identifiers must start with either a letter or an underscore but can afterwards be followed by digits.
When handling identifiers we start off by consuming the entire literal.
This is in accordance with the maximal munch principal where a token is the longest possible valid string.
The entire character sequence is matched against the map of all reserved keywords.
Only if the sequence does not match any reserved keyword it is categorized as an identifier.
This order of operations is deliberate to ensure that keywords have a higher priority than user-defined identifiers.

%! Kan bare slette dette også. Prøvede at lave en nice overgang, men måske er det ikke nice.
If no errors are detected, we know the input was without lexical errors and is now converted into a 
sequence of tokens that we can use in the next step: Parsing.